\documentclass[12pt, aspectratio=169]{beamer}

\input{../header}
% \usepackage{booktabs}
% \usepackage{tikz}
% \usetikzlibrary{patterns}
% \DeclareMathOperator{\E}{E}
% \DeclareMathOperator{\Var}{Var}
% \DeclareMathOperator{\V}{V}
% \DeclareMathOperator{\C}{C}


\title{Sampling Methods}
\author{Md. Aminul Islam Shazid}
\date{}


\begin{document}
    {
		\setbeamertemplate{footline}{}    % NO FOOTLINE FOR THESE TWO FRAMES
		\addtocounter{framenumber}{-2}    % not counting the title page and the outline in frame numbers

		\begin{frame}
			\titlepage
		\end{frame}

		\begin{frame}{Outline}
            \vfill
			\tableofcontents[subsectionstyle=hide]
            \vfill
		\end{frame}
	}


	\section{Introduction}
    \begin{frame}{Sampling}
        \begin{itemize}
            \item In many practical situations it is difficult or impossible to observe every unit in a population

            \item Instead, we select a subset of the population and collect information from that subset

            \item This process is called \textbf{sampling}

            \item The selected subset of the population is called a \textbf{sample}
            \end{itemize}

        \begin{block}{Definition}
            Sampling is the process of selecting a representative subset of units from a population in order to study the characteristics of the population
        \end{block}
    \end{frame}


    \begin{frame}{Basic Concepts in Sampling}
        \begin{itemize}
            \item \textbf{Population}

            The entire collection of individuals or units about which information is desired.

            \item \textbf{Sample}

            A subset of the population selected for analysis

            \item \textbf{Sampling Method}

            The procedure used to select the sample from the population
        \end{itemize}
    \end{frame}


    \begin{frame}{Additional Sampling Terminology}
        \begin{itemize}
            \item \textbf{Sampling Unit}

            The basic element or group of elements considered for selection.

            Example: households, students, firms.

            \item \textbf{Sampling Frame}

            A complete list of all sampling units in the population.

            \item \textbf{Sample Unit}

            A sampling unit that is actually selected into the sample.

            \item \textbf{Unit of Inquiry}

            The unit about which information is collected. It may differ from the sampling unit.
        \end{itemize}
    \end{frame}


    \begin{frame}{Sampling Methods}
        Sampling methods can broadly be classified into two categories:

        \begin{itemize}
            \item \textbf{Probability Sampling}

            Each unit in the population has a known and non-zero probability of being selected.

            \item \textbf{Non-probability Sampling}

            Selection of units depends on judgement, convenience, or other non-random mechanisms.
        \end{itemize}
    \end{frame}


    \section{Probability Sampling}

    \begin{frame}{Probability Sampling}
        \begin{itemize}
            \item Each element of the population has a known probability of being selected

            \item Randomization is used to select sample units

            \item Sampling bias can be reduced

            \item Statistical inference about the population becomes possible
        \end{itemize}

        \textbf{Common probability sampling methods}

        \begin{itemize}
            \item Simple Random Sampling
            \item Stratified Random Sampling
            \item Systematic Sampling
            \item Cluster Sampling
        \end{itemize}
    \end{frame}


    \subsection{Simple Random Sampling}

    \begin{frame}{Simple Random Sampling (SRS)}
        \begin{block}{Definition}
            A sampling method in which every unit in the population has an equal probability of being selected.
        \end{block}

        \begin{itemize}
            \item If the population size is $N$, the probability of selecting any specific unit is
            \[
                \frac{1}{N}
            \]

            \item Every possible sample of size $n$ has equal probability of being selected
        \end{itemize}
    \end{frame}


    \begin{frame}{Common Selection Methods in SRS}
        Selection method means how items are selected into the sample.
        \begin{itemize}
            \item Lottery method
            \item Random number table
            \item Computer generated random numbers
        \end{itemize}
    \end{frame}


    \subsection{Stratified Random Sampling}

    \begin{frame}{Stratified Random Sampling}
        \begin{block}{Definition}
            In stratified sampling the population is divided into several non-overlapping groups called \textbf{strata}, and a random sample is selected from each stratum.
        \end{block}

        \begin{itemize}
            \item Strata are formed using known characteristics

            \item Units within each stratum are relatively homogeneous
            
            \item Units across different strata ar heterogeneous

            \item Sampling is performed independently within each stratum
        \end{itemize}
    \end{frame}


    \begin{frame}{Procedure of Stratified Random Sampling}
        \begin{enumerate}
            \item Divide the population into $L$ strata

            \item Determine total sample size $n$

            \item Allocate sample sizes $n_1, n_2, \dots, n_L$ to each stratum

            \item Select a random sample from each stratum
        \end{enumerate}

        \begin{align*}
            \text{Population sizes of each strata: } & N_1 + N_2 + \cdots + N_L = N \\
            \text{Sample sizes of each strata: } & n_1 + n_2 + \cdots + n_L = n
        \end{align*}
    \end{frame}


    \subsection{Systematic Sampling}

    \begin{frame}{Systematic Sampling}
        \begin{block}{Idea}
            Sample units are selected at regular intervals from an ordered list of the population.
        \end{block}

        \begin{itemize}
            \item Population size: $N$

            \item Sample size: $n$

            \item Sampling interval:
            \[
                k = \frac{N}{n}
            \]
        \end{itemize}
    \end{frame}


    \begin{frame}{Procedur of Systematic Sampling}
        \begin{enumerate}
            \item Number the population units from $1$ to $N$

            \item Compute the sampling interval $k = N/n$

            \item Randomly select the first unit from $1,2,\dots,k$

            \item Select subsequent units as
            \[
                r, r+k, r+2k, \dots
            \]
        \end{enumerate}
    \end{frame}


    \subsection{Cluster Sampling}

    \begin{frame}{Cluster Sampling}
        \begin{block}{Definition}
            Cluster sampling is a method in which the population is divided into groups called clusters and a random sample of clusters is selected.
        \end{block}

        \begin{itemize}
            \item Each cluster contains multiple population elements

            \item All elements in selected clusters may be included in the sample

            \item Often used when a complete sampling frame is not available
        \end{itemize}
    \end{frame}


    \section{Non-probability Sampling}

    \begin{frame}{Non-probability Sampling}
        \begin{itemize}
            \item In non-probability sampling, the probability of selection is not known

            \item Selection depends on judgement, accessibility, or convenience

            \item These methods are easier to implement but may introduce bias
        \end{itemize}

        \textbf{Common types}

        \begin{itemize}
            \item Convenience sampling
            \item Purposive sampling
            \item Quota sampling
            \item Snowball sampling
        \end{itemize}
    \end{frame}


    \begin{frame}{Examples of Non-probability Sampling}
        \begin{itemize}
            \item \textbf{Convenience Sampling}

            Selecting individuals who are easiest to reach.

            \item \textbf{Purposive Sampling}

            Selecting individuals based on specific characteristics.

            \item \textbf{Quota Sampling}

            Selecting units to satisfy predefined group proportions.

            \item \textbf{Snowball Sampling}

            Existing participants recruit additional participants.
        \end{itemize}
    \end{frame}


    \begin{frame}{Summary}
        \begin{itemize}
            \item Sampling allows us to study large populations efficiently

            \item Sampling methods are broadly classified as probability and non-probability methods

            \item Probability sampling supports statistical inference

            \item Common probability methods include simple random, stratified, systematic, and cluster sampling
        \end{itemize}
    \end{frame}


    \section*{Thank you.}

\end{document}